\section{Methodology}
% Must also include:
% Feature Selection: how features are chosen? in case a subset is chosen to
% solve your task describe the model selection strategy
\begin{itemize}
    \item What do I try to maximize (or better on what do I wanto to improve
    upon? e.g. f2 score of the baseline model.)
    \item How am I doing it? (one example: running a cnn, when not sure about the class, run a expert that destinguishes between those two classes)
\end{itemize}
I think in this chapter I really want to tell a story. For that I will start
with the baseline model and take a look at the confusion matrix to see where the
model is struggling. Then I will try to find a way to improve the model on those
classes with techniques that I have learned in the course and techniques that I
have found in the literature. Especially going to try to improve on f2 on the
cancer class itself but also try to look for some correlation between the
classes themselves. Furthermore, if the model is not sure about a class, e.g.
having a confidence < 0.9 then a second model will be used which is an expert on
distinguishing between those two classes. This model will use the intermediate
layers of the baseline model.

I guess even further after training the baseline model and seeing where it
struggles I have to talk about interpretability and try to find out why the
model is struggling on those classes.

When I use literature I need to cite the respective paper.
